\documentclass{article}
\usepackage[UTF8]{ctex}
\usepackage{amsmath,amssymb,amsfonts}
\usepackage{enumitem}
\usepackage{geometry}
\geometry{a4paper, margin=1in}
% === 错题本排版优化（不参与题目标识） ===
\setlength{\parindent}{0pt}
\setlength{\parskip}{0.35em plus 0.10em minus 0.08em}
\setlength{\emergencystretch}{3em}
\setlist{itemsep=0.25em, topsep=0.40em, parsep=0pt}
\allowdisplaybreaks[2]
% === 错题本排版优化结束 ===
\begin{document}
\section*{南京航空航天大学2024年数学分析真题}
\section*{98．南京航空航天大学2024年数学分析试卷}

\subsection{题目1}
1．解答如下问题：

\begin{itemize}
  \item[（1）] 求极限 \(\lim _{x \rightarrow 0} \frac{1}{x^{4}}\left[6+x^{2}-\frac{5}{6} x^{4}-6 \sqrt[3]{2-\cos x}\right]\) 。
  \item[（2）] 求定积分 \(\int_{0}^{\frac{\pi}{4}} \frac{x}{1+\cos 2 x} \mathrm{~d} x\) 。
\end{itemize}

\subsection{题目2}
2．设 \(f(x)=\frac{x+2}{x+1} \sin \frac{1}{x}, a>0\) 为常数，证明：\(f(x)\) 在 \([a,+\infty)\) 上一致连续，在 \((0, a)\) 上不一致连续．

\begin{itemize}

\subsection{题目3}
\end{itemize}

\subsection{题目3}
原题边界“3．”：设 \(f(x)\) 在 \([a, b]\) 上连续，在 \((a, b)\) 内二阶可导，证明：对任意的 \(x \in[a, b]\) ，存在 \(\xi \in[a, b]\) ，使得
\[
\frac{f(x)-f(a)}{x-a}-\frac{f(b)-f(a)}{b-a}=\frac{1}{2} f^{\prime \prime}(\xi)(x-b) .
\]

\subsection{题目4}

\subsection{题目4}
原题边界“4．”：判断反常积分 \(\int_{0}^{+\infty} \frac{\left(\ln ^{2} x\right) \sin x}{x+1} \mathrm{~d} x\) 是绝对收敛，条件收敛还是发散，并证明．

\subsection{题目5}
5．解答如下问题：

\begin{itemize}
  \item[（1）] 设 \(f(x), g(x)\) 在 \([a, b]\) 上可积，证明 Cauchy－Schwarz 不等式：
\[
\left(\int_{a}^{b} f(x) g(x) \mathrm{d} x\right)^{2} \leq \int_{a}^{b} f^{2}(x) \mathrm{d} x \int_{a}^{b} g^{2}(x) \mathrm{d} x
\]
  \item[（2）] 设 \(f(x)\) 在 \([a, b]\) 上有连续的导函数，且 \(f(a)=0\) ，证明：
\[
\int_{a}^{b}\left|f(x) f^{\prime}(x)\right| \mathrm{d} x \leq \frac{b-a}{2} \int_{a}^{b}\left[f^{\prime}(x)\right]^{2} \mathrm{~d} x .
\]

\subsection{题目6}
\end{itemize}

\subsection{题目6}
原题边界“6．”：设 \(a_{1}=2, a_{n+1}=\frac{1}{2}\left(a_{n}+\frac{1}{a_{n}}\right), n=1,2, \cdots\) ．

\begin{itemize}
\item[（1）] 证明 \(\left\{a_{n}\right\}\) 收敛，并求其极限．
  \item[（2）] 判断级数 \(\sum_{n=1}^{\infty}\left(\frac{a_{n}}{a_{n+1}}-1\right)\) 的敛散性并证明．
\end{itemize}

\subsection{题目7}
7．将 \(f(x)=e^{x}(-\pi \leq x \leq \pi)\) 展开成傅里叶级数，讨论级数的收敛性，并求级数 \(\sum_{n=0}^{\infty} \frac{1}{1+n^{2}}\) 的和．

\subsection{题目8}
8．给定方程组

\[
\left\{\begin{array}{l}
x u^{2}+v=y^{3} \\
2 y u-x v^{3}=4 x
\end{array}\right.
\]

判断其在 \(P_{0}\left(x_{0}, y_{0}, u_{0}, v_{0}\right)=(0,1,0,1)\) 附近能否确定 \(u, v\) 为 \(x, y\) 的函数，若能，求出偏导数\(\frac{\partial u}{\partial y}, \frac{\partial v}{\partial y}, \frac{\partial^{2} v}{\partial y^{2}}\) 在点 \(\left(x_{0}, y_{0}\right)=(0,1)\) 的值．

\subsection{题目9}
9．解答如下问题：

\begin{itemize}
  \item[（1）] 求曲面积分
\[
\iint_{S}\left(x^{2}+y^{2}\right) \mathrm{d} S
\]
其中 \(S\) 为立体 \(\sqrt{x^{2}+y^{2}} \leq z \leq 1\) 的边界曲面。
  \item[（2）] 求曲线积分
\[
\oint_{L}(y+\sin x) \mathrm{d} x+\left(z-e^{y}\right) \mathrm{d} y+(x+1) \mathrm{d} z .
\]
其中 \(L\) 为球面 \(x^{2}+y^{2}+z^{2}=R^{2}\) 与平面 \(x+y+z=0\) 的交线，从 \(x\) 轴正向看去，\(L\) 沿逆时针方向。

\subsection{题目10}
\end{itemize}

\subsection{题目10}
原题边界“10．”：设 \(f(x)\) 在 \([0,+\infty)\) 上可导，且 \(f(0)=0\) ．求极限
\[
\lim _{t \rightarrow 0^{+}} \frac{1}{\pi t^{4}} \iiint_{V} f\left(\sqrt{x^{2}+y^{2}+z^{2}}\right) \mathrm{d} x \mathrm{~d} y \mathrm{~d} z .
\]
其中 \(V: x^{2}+y^{2}+z^{2} \leq t^{2}\) ．

\subsection{题目11}

\subsection{题目11}
原题边界“11．”：设 \(\left\{f_{n}(x)\right\}\) 是在 \([a, b]\) 上连续的函数列，且 \(\left\{f_{n}(x)\right\}\) 在 \([a, b]\) 上一致收敛于 \(f(x), f(x)\) 在 \([a, b]\) 上恒为负值．证明：当 \(n\) 充分大时，\(f_{n}(x)\) 在 \([a, b]\) 上也恒为负值，且函数列 \(\left\{\frac{1}{f_{n}(x)}\right\}\) 在 \([a, b]\) 上一致收敛于 \(\frac{1}{f(x)}\) ．

\subsection{题目12}

\subsection{题目12}
原题边界“12．”：设 \(f(x)\) 在 \([a, b]\) 上可积，且 \(A \leq f(x) \leq B, g(u)\) 在 \([A, B]\) 上连续，证明：\(g(f(x))\) 在 \([a, b]\) 上可积．
\end{document}
